% year: תשע"ט
% semester: א
% moed: א
% number: 3ב
% points: 10
% categories: func-limits, sequences
% source: exams/מבחנים/תשעט/מועד א תשעט.pdf

\begin{question}
(10 נק') הסבירו מדוע לפונקציה $f(x)=\cos\left(\dfrac{x+1}{x-1}\right)$ אין גבול כאשר $x$ שואף ל-$1$.
\end{question}

\begin{hint}
השתמשו בהגדרת הגבול לפי היינה: מצאו שתי סדרות השואפות ל-$1$ שעליהן ערכי הפונקציה שואפים לגבולות שונים.
\end{hint}

\begin{solution}
נשתמש בהגדרת היינה. אם היה קיים $L=\lim_{x\to1}f(x)$, אז לכל סדרה $x_n\to1$, $x_n\ne1$, היה מתקיים $f(x_n)\to L$.

נשים לב שאם $\frac{x+1}{x-1}=y$ אז $x=\frac{y+1}{y-1}$. נבחר
\[ x_n=\frac{2\pi n+1}{2\pi n-1},\qquad z_n=\frac{(2n+1)\pi+1}{(2n+1)\pi-1}. \]
שתי הסדרות שואפות ל-$1$ (מחלקים מונה ומכנה ב-$n$) ושונות מ-$1$. מתקיים $\frac{x_n+1}{x_n-1}=2\pi n$ ו-$\frac{z_n+1}{z_n-1}=(2n+1)\pi$, ולכן
\[ f(x_n)=\cos(2\pi n)=1\to1,\qquad f(z_n)=\cos\big((2n+1)\pi\big)=-1\to-1. \]
קיבלנו שתי סדרות השואפות ל-$1$ שעליהן ערכי הפונקציה מתכנסים לגבולות שונים, בסתירה ליחידות הגבול. לכן ל-$f$ אין גבול ב-$1$.

(באופן אינטואיטיבי: כאשר $x\to1^\pm$ הביטוי $\frac{x+1}{x-1}\to\pm\infty$, והקוסינוס מתנודד בין $-1$ ל-$1$ אינסוף פעמים.)
\end{solution}
